% !TEX root = ../main.tex
% =====================================================================
% Section 2 — Related Work  (gdn2vessel, benchmark-only D&B 稿)
% 重定位 2026-06-24：benchmark-only（STORY_FRAMEWORK.md 真源，两命门全死）。
% 框架 = 支撑「为何需要这个 benchmark」，每条点名划界；禁卖方法/机制。
% 所有 SOTA 数字 = \TODO 占位（禁借数当我们的，禁臆造）；citation key 占位。
% 双盲：无作者/机构/可识别自引。
% 防御规则 R1/R2/R3/R3b/R4/R6（见 STORY_FRAMEWORK.md）：
%   - novelty 收窄：禁 "first disconnection benchmark"（PTR 已占 3D 肺）；
%     禁 "first in-model connectivity"（CorSegRec/GLCP 已做 reconnection）。
%   - 记忆模块/Frangi 门 = leaderboard 被评测项，不卖机制特异（已证伪）。
%   - 区分度 claim 一律 hedge（批2 未跑）：we report what the benchmark reveals。
% =====================================================================

\section{Related Work}
\label{sec:related}

We position our benchmark against five threads of prior work. The first two
threads---vessel/tubular segmentation backbones and topology- or
connectivity-oriented methods---supply the \emph{methods that our leaderboard
evaluates}; they are subjects of measurement, not competing benchmarks. The
remaining three threads---prior synthetic-disconnection benchmarks, the closest
2D fundus reconnection method, and benchmark/metric-diagnosis studies---are where
our contribution must be \emph{delineated}, and we do so explicitly below. We
make no claim that any single component (a segmentation backbone, a connectivity
loss, or a memory module) is novel in isolation; the contribution we defend is
the \emph{evaluation suite}---a reproducible synthetic-break protocol, a family
of reconnection metrics with a diagnosed failure mode, a full-spectrum
leaderboard, and a pre-registered diagnostic methodology---specialised to the
\emph{2D fundus} setting.

% ---------------------------------------------------------------------
\subsection{Vessel and Tubular Segmentation Backbones}
\label{sec:related:vessel}

\paragraph{CNN backbones.}
Curvilinear and tubular segmentation has long been driven by encoder--decoder
CNNs. Beyond U-Net and its variants, specialised designs target the thin,
elongated geometry of vessels: IterNet~\cite{key:iternet}, CS$^2$-Net~\cite{key:csnet},
and FR-UNet~\cite{key:frunet} refine multi-scale context, attention, or iterative
refinement. These networks optimise per-pixel delineation and form a part of the
classical-CNN tier of our leaderboard.

\paragraph{State-space (Mamba) backbones.}
A more recent line replaces or augments convolutional context with selective
state-space models (SSMs). Generic medical backbones include
VM-UNet~\cite{key:vmunet}, U-Mamba~\cite{key:umamba}, and
Swin-UMamba~\cite{key:swinumamba}; vessel- and tubular-specific variants such as
HM-Mamba~\cite{key:hmmamba} couple hierarchical multi-scale tubular convolution
with frequency-domain gating, and related scan-ordering designs emphasise the
spatial traversal itself. We treat all of these as \emph{baselines to be
evaluated} rather than as competitors to our benchmark: they define the
SSM/Mamba and CNN tiers of the full-spectrum leaderboard
(Section~\ref{sec:leaderboard}). We report each method's reproduced segmentation
scores against published values for sanity (any reproduction gap of more than a
few points is flagged, not silently claimed as a win), and place every method on
the same disconnection--reconnection protocol so that they are scored on
identical synthetic breaks.

% ---------------------------------------------------------------------
\subsection{Topology, Connectivity, and Vessel Reconnection}
\label{sec:related:topo}

\paragraph{Connectivity losses and topological objectives.}
A complementary body of work enforces connectivity through dedicated objectives
rather than per-pixel overlap. clDice~\cite{key:cldice} and its radius-aware
extension cbDice~\cite{key:cbdice} provide differentiable soft-skeleton losses;
persistent-homology and Betti-matching objectives~\cite{key:bettimatching}
penalise topological discrepancies; and Skeleton Recall~\cite{key:skeletonrecall}
offers a lightweight skeleton objective. We do not propose a new connectivity
loss. Instead, clDice, Betti errors, and Skeleton Recall enter our metric family
as standard topology axes that cross-check our reconnection-specific measures
(Section~\ref{sec:bench}).

\paragraph{Vessel reconnection methods (delineation from our scope).}
Crucially, reconnecting broken vessels is \emph{not} an unaddressed problem, and
we therefore make \textbf{no} claim of being the first to perform in-model
connectivity or vessel reconnection. CorSegRec~\cite{key:corsegrec} couples
segmentation with reconnection and introduces its own continuity-oriented score,
and GLCP~\cite{key:glcp} (MICCAI~2025 Oral) explicitly models \emph{local
breakpoints together with global connectivity}. These are mature reconnection
methods; in our work they are candidate entries on the leaderboard rather than
points of novelty. Our distinction is one of \emph{evaluation rather than
method}: where these works contribute reconnection \emph{models}, we contribute a
standardised suite---a controlled synthetic-break protocol with a severity grid,
a reconnection metric family with a diagnosed failure mode, and a leaderboard
that places such methods on common ground.

% ---------------------------------------------------------------------
\subsection{Synthetic-Disconnection Benchmarks and Tracing}
\label{sec:related:bench}

\paragraph{Prior disconnection benchmarks (the precedent we narrow against).}
The idea of a benchmark built on \emph{synthetic} structural breaks already
exists. PTR~\cite{key:ptr} (MICCAI~2023), released as ``baselines and a
dataset,'' provides a synthetic-disconnection benchmark for \emph{3D pulmonary}
tubular structures. Because this precedent occupies the disconnection-benchmark
form, we deliberately \textbf{do not} claim ``the first disconnection
benchmark.'' We narrow our headline to the first \emph{2D fundus} reconnection
\emph{evaluation suite}, defined by the specific combination of a reproducible
synthetic-break protocol, a reconnection metric family (with metric-failure
diagnosis), a full-spectrum baseline leaderboard, and a pre-registered
diagnostic methodology---a combination, to our knowledge, not previously
assembled for retinal fundus images.

\paragraph{Tracing-based connectivity.}
A separate family recovers connectivity by \emph{tracing} curvilinear structures,
e.g.\ graph-based vessel networks~\cite{key:vgn} and transformer tracing models
such as Trexplorer~\cite{key:trexplorer}. These methods typically require seed
points or a starting node and produce graphs rather than dense masks; they
address a related but distinct task and are not the object of our mask-based
reconnection protocol.

% ---------------------------------------------------------------------
\subsection{2D Fundus Completion and Our Protocol Source}
\label{sec:related:fundus}

\paragraph{Closest 2D fundus reconnection method.}
The closest neighbour in modality and task is MaskVSC~\cite{key:maskvsc}
(TMI~2025), which performs masked completion of 2D retinal vessels. We delineate
our work from it on the same method-versus-evaluation axis: MaskVSC is a
\emph{method paper} that proposes a completion model, whereas we contribute an
\emph{evaluation suite}---a protocol, a metric family with a documented failure
mode, a leaderboard, and a diagnostic methodology---on which methods of this kind
can be measured.

\paragraph{Source of our synthetic-break protocol.}
Our disconnection protocol follows the plug-and-play synthetic-break generation
of Corvo~\emph{et al.}~\cite{key:creatis} (arXiv~2404.10506). We stress that this
is a \emph{reference source for our protocol, not a competitor}: their work is a
method paper, and we adapt its break-generation procedure (with the
radius-distribution and gap parameters fixed for reproducibility, organised into
a four-level severity grid) as the controlled stimulus of our benchmark. We note
two deviations from the original protocol---it includes neither a boundary-blur
step nor a success-rate metric---so the success-rate metric we report is our own
addition, defined explicitly and accompanied by its failure-mode diagnosis
(Section~\ref{sec:bench}). The protocol and all evaluated methods derive breaks
and predictions from the input image only, never from ground-truth topology.

% ---------------------------------------------------------------------
\subsection{Benchmark and Metric-Diagnosis Studies}
\label{sec:related:diag}

A growing line of dataset-and-benchmark work shows that \emph{systematic
evaluation revealing the limits of methods and metrics} is itself a contribution.
Touchstone~\cite{key:touchstone} (NeurIPS~2024 Datasets \& Benchmarks) re-ranks
segmentation methods under a controlled large-scale protocol, and recent
pixel-wise-metric analyses~\cite{key:pixelwise} (NeurIPS~2025) expose how
overlap-style metrics can mislead. These precedents legitimise our two
diagnostic contributions: (i) the finding that our success-rate metric can be
\emph{saturated by over-prediction}---a degenerate, low-overlap model can score
near the ceiling because the gap-closure check only inspects coverage of the
erased region---which is why we use a topological-continuity measure as the
primary reconnection axis and report success rate only alongside a specificity or
over-prediction context; and (ii) our pre-registered honest negative results
(Section~\ref{sec:diag}). Following the same metric-diagnosis tradition, we treat
these findings as evidence about \emph{what the benchmark reveals} rather than as
claims that any one method or metric is superior.

% ---------------------------------------------------------------------
\subsection{Linear-Attention Memory Modules (an Evaluated Method, Not a Claim)}
\label{sec:related:memory}

For completeness we situate the memory-based entry on our leaderboard. Linear
attention and delta-rule sequence models replace softmax attention with a
recurrent matrix-valued state updated in linear time, including gated linear
attention (GLA)~\cite{key:gla} and the Gated DeltaNet family. The entry we label
\texttt{ours\_gdn2} is built on GDN-2~\cite{key:gdn2}, a gated delta-rule
associative memory; a separate temporal use of gated delta KV memory appears in
GDKVM~\cite{key:gdkvm} (ICCV~2025), which indexes memory across echocardiography
\emph{video frames}, whereas our single-image setting indexes a \emph{spatial}
scan and has no temporal axis. We emphasise---consistent with our honest
reporting (Section~\ref{sec:diag})---that \texttt{ours\_gdn2} is simply
\emph{one evaluated method on the leaderboard}, and we make \textbf{no} claim of
mechanism specificity for the delta-rule memory. On the contrary, our
pre-registered audit finds that delta-rule memory is \emph{not} specific on 2D
vessels: at the only comparable convergence point, gated delta memory matches
plain gated linear attention on an associative-recall probe, and the
segmentation objective provides no gradient incentive that would bind the memory
state to vessel identity. This negative result is reported as evidence that the
benchmark has discriminative power, not as a competitive advantage.
